\chapter{Geometric structures on manifolds}

\begin{definition}[Pseudogroup of local models]
\label{thurston-ch03-pseudogroup}
\group{geometric-structures}
\source{thurston:page-0037}
\notready
A pseudogroup \(\G\) on a manifold \(X\) is a collection of local homeomorphisms between open subsets of \(X\), closed under restriction, composition where defined, inversion, and gluing over open covers. A \(\G\)-manifold is a manifold obtained by gluing open subsets of \(X\) with transition maps in \(\G\).
\end{definition}

\begin{definition}[Homogeneous geometric structure]
\label{thurston-ch03-gx-manifold}
\group{geometric-structures}
\source{thurston:page-0037}
\notready
If a group \(G\) acts on \(X\), a \((G,X)\)-manifold is a \(\G\)-manifold for the pseudogroup of restrictions of elements of \(G\) to open subsets of \(X\).
\end{definition}

\begin{remark}[Basic model geometries]
\label{thurston-ch03-basic-models}
\group{geometric-structures}
\source{thurston:page-0037,thurston:page-0038,thurston:page-0039}
\notready
Local \(C^r\)-diffeomorphisms, piecewise-linear homeomorphisms, affine transformations, Euclidean isometries, orthogonal transformations of elliptic space, and hyperbolic isometries give respectively differentiable, PL, affine, Euclidean, elliptic, and hyperbolic manifolds. Every closed Euclidean manifold is finitely covered by a torus; a closed Euclidean manifold is determined up to affine equivalence by its fundamental group. Every surface of negative Euler characteristic admits a hyperbolic structure.
\end{remark}

\begin{conjecture}[Finite fundamental group conjecture]
\label{thurston-ch03-elliptic-conjecture}
\group{geometric-structures}
\source{thurston:page-0038,thurston:page-0039}
\notready
Every three-manifold with finite fundamental group admits an elliptic structure.
\end{conjecture}

\begin{proposition}[Ideal tetrahedral figure-eight construction]
\label{thurston-ch03-figure-eight-tetrahedra}
\dcref{3.1}
\group{hyperbolic-constructions}
\source{thurston:page-0039,thurston:page-0040,thurston:page-0041}
\notready
The ideal hyperbolic tetrahedron obtained from a regular Euclidean tetrahedron inscribed in the sphere at infinity has all dihedral angles \(\pi/3\). Two copies can be glued along their faces in the standard figure-eight pattern, producing a hyperbolic structure on the figure-eight knot complement; its fundamental group is a subgroup of index \(12\) in \(\PSL_2(\Z[\omega])\), where \(\omega\) is a primitive cube root of unity.
\end{proposition}
\begin{proof}
\sketch The bounding circles of the extended faces meet at \(\pi/3\), hence the hyperbolic dihedral angles have that value. The face pairings are hyperbolic isometries, and the six incident angles around each edge sum to \(2\pi\), so the quotient is locally hyperbolic. The stated arithmetic identification is the known group description of the resulting quotient.
\end{proof}

\begin{proposition}[Geodesic-boundary tetrahedral construction]
\label{thurston-ch03-geodesic-boundary}
\dcref{3.2}
\group{hyperbolic-constructions}
\source{thurston:page-0041,thurston:page-0042}
\notready
There is a hyperbolic three-manifold \(M\), with genus-two totally geodesic boundary, obtained by gluing two truncated regular tetrahedra whose original dihedral angles are \(\pi/6\) and whose truncation faces meet the remaining faces orthogonally. Doubling \(M\) along its boundary gives a closed hyperbolic three-manifold.
\end{proposition}
\begin{proof}
\sketch In the projective model, expand a regular ideal tetrahedron until its edges are tangent to the sphere at infinity; continuity changes the dihedral angles from \(\pi/3\) to \(0\), so an intermediate tetrahedron has angle \(\pi/6\). Truncation at the four polar planes creates right angles, and the prescribed gluing makes the interior edge angles equal \(2\pi\). The truncation faces form the totally geodesic boundary, whose double is closed.
\end{proof}

\begin{proposition}[Octahedral link-complement constructions]
\label{thurston-ch03-octahedral-links}
\dcref{3.3,3.4}
\group{hyperbolic-constructions}
\source{thurston:page-0042,thurston:page-0043,thurston:page-0044}
\notready
The Whitehead link complement admits a hyperbolic structure from one ideal regular octahedron, whose dihedral angles are \(\pi/2\), with faces paired in the indicated pattern. The Borromean rings complement admits one from two ideal regular octahedra, paired by corresponding faces with alternating rotations of \(2\pi/3\).
\end{proposition}
\begin{proof}
\sketch A regular Euclidean octahedron inscribed in the sphere becomes an ideal hyperbolic octahedron with right dihedral angles. The displayed face pairings are hyperbolic isometries; their edge angle sums give the stated link complements. The same local angle check applies to the two-octahedron Borromean gluing, where the rotations alternate consistently around the complex.
\end{proof}

\begin{theorem}[Developing map and holonomy]
\label{thurston-ch03-developing-map}
\dcref{3.5}
\group{developing-holonomy}
\source{thurston:page-0044,thurston:page-0045}
\notready
Let \(X\) be real analytic and let \(G\) be a group of real analytic diffeomorphisms of \(X\). Every \((G,X)\)-structure on a connected manifold \(M\) has a developing map \(D:\widetilde M\to X\), obtained by analytic continuation of a chart, and \(D\) is a local \((G,X)\)-homeomorphism, unique up to postcomposition by an element of \(G\). Each deck transformation \(T_\alpha\) satisfies \(D\circ T_\alpha=g_\alpha\circ D\), and \(H(\alpha)=g_\alpha\) defines a homomorphism \(H:\pi_1(M)\to G\), the holonomy.
\end{theorem}
\begin{proof}
\sketch On overlaps, analytic transition maps are restrictions of elements of \(G\), and analyticity makes each continuation element locally constant. Continue a chart along paths and pass to path classes to obtain \(D\) on \(\widetilde M\). Uniqueness of continuation gives the equivariance equation; composing deck transformations proves \(g_{\alpha\beta}=g_\alpha g_\beta\).
\end{proof}

\begin{definition}[Complete geometric structure]
\label{thurston-ch03-complete-structure}
\group{developing-holonomy}
\source{thurston:page-0045}
\notready
A \((G,X)\)-manifold is complete when its developing map \(D:\widetilde M\to X\) is a covering map. If \(X\) is simply connected, completeness means that \(D\) is a homeomorphism.
\end{definition}

\begin{proposition}[Reconstruction from complete holonomy]
\label{thurston-ch03-complete-quotient}
\group{developing-holonomy}
\source{thurston:page-0045}
\notready
If \(X\) is simply connected and a complete \((G,X)\)-manifold has holonomy image \(\Gamma\), then it is reconstructed as the quotient \(X/\Gamma\) (with the induced action).
\end{proposition}
\begin{proof}
\sketch The developing map identifies the universal cover with \(X\); deck equivariance identifies the deck group with \(\Gamma\), so passing to the quotient gives \(M\cong X/\Gamma\).
\end{proof}

\begin{proposition}[Compact transitive structures are complete]
\label{thurston-ch03-compact-isotropy-complete}
\dcref{3.6}
\group{completeness}
\source{thurston:page-0045,thurston:page-0046}
\notready
Let \(G\) act transitively and analytically on \(X\), with compact isotropy groups. Every closed \((G,X)\)-manifold is complete.
\end{proposition}
\begin{proof}
\sketch Average a positive definite form at one point over its compact isotropy group, then transport it by transitivity to obtain a \(G\)-invariant Riemannian metric on \(X\), hence on \(M\) and \(\widetilde M\). A sufficiently small metric ball in closed \(M\) is uniformly convex. Its inverse image under \(D\) is a disjoint union of copies of a half-sized ball; path lifting then extends across every endpoint, so \(D\) is a covering.
\end{proof}

\begin{proposition}[Metric criteria for completeness]
\label{thurston-ch03-metric-completeness}
\dcref{3.7}
\group{completeness}
\source{thurston:page-0046,thurston:page-0048}
\notready
For a \((G,X)\)-manifold under the hypotheses of Proposition~\ref{thurston-ch03-compact-isotropy-complete}, the following are equivalent: (a) the natural metric on \(M\) is complete; (b) some \(\epsilon>0\) has all closed \(\epsilon\)-balls compact; (c) every closed \(k\)-ball is compact for every \(k>0\); (d) there are compact sets \(S_t\) exhausting \(M\) such that \(S_{t+a}\) contains the radius-\(a\) neighborhood of \(S_t\).
\end{proposition}
\begin{proof}
\sketch Completeness of the metric on \(\widetilde M\) lets one lift any path in \(X\): the set of times for which a lift exists is open by local invertibility and closed by metric completeness. Thus the developing map covers. The implications among (b)--(d) are elementary, (d) implies (a), and transitivity supplies a uniform compact ball in the model. A complete quotient of the model inherits (b), proving equivalence.
\end{proof}

\begin{definition}[Horosphere and horoball]
\label{thurston-ch03-horosphere}
\dcref{3.8}
\group{hyperbolic-geometry}
\source{thurston:page-0048,thurston:page-0049}
\notready
A horosphere centered at a point \(X\) at infinity in hyperbolic space is a hypersurface orthogonal to every geodesic ending at \(X\); its bounded region is a horoball. In the upper half-space model, horospheres centered at infinity are horizontal Euclidean planes, and each horosphere is isometric to \(E^{n-1}\). Isometries fixing the center act on a horosphere by similarities.
\end{definition}

\begin{lemma}[Horospherical scaling]
\label{thurston-ch03-horosphere-scaling}
\group{hyperbolic-geometry}
\source{thurston:page-0049,thurston:page-0050}
\notready
If concentric horospheres \(h_1,h_2\) are hyperbolic distance \(d\) apart, then distances between corresponding orthogonal geodesics measured on \(h_2\) and \(h_1\) have ratio \(e^d\) (with the upper horosphere in the numerator).
\end{lemma}
\begin{proof}
\sketch In upper half-space, the metric is \(ds^2=dx^2/x_n^2\). Horizontal lengths scale inversely with height, while vertical distance between heights \(h_1,h_2\) is \(\abs{\log h_2-\log h_1}\), giving the exponential ratio.
\end{proof}

\begin{theorem}[Completeness of ideal-triangle surfaces]
\label{thurston-ch03-ideal-triangle-completeness}
\dcref{3.9}
\group{hyperbolic-surfaces}
\source{thurston:page-0050,thurston:page-0051,thurston:page-0052}
\notready
All ideal hyperbolic triangles are congruent. For an ideal vertex \(v\) in a surface \(S\) assembled from ideal triangles, let \(d(v)\) be the signed distance between an initial horocycle and its concentric return after one circuit around \(v\), positive when the returned horoball contains the initial one. Then \(S\) is complete if and only if \(d(v)=0\) for every ideal vertex. If some \(d(v)\ne0\), the metric completion has one geodesic boundary component of length \(\abs{d(v)}\) for each such vertex.
\end{theorem}
\begin{proof}
\sketch If \(d(v)<0\), successive circuits contract horocycle lengths by a fixed factor, producing a bounded nonconvergent Cauchy sequence. The case \(d(v)>0\) is detected after reversing the circuit. If all invariants vanish, truncate all ends by horoballs and expand the truncation uniformly; the resulting compact exhaustion satisfies condition (d) of Proposition~\ref{thurston-ch03-metric-completeness}. Nonzero invariants produce spiralling horocycles whose metric completion adds the stated boundary circles.
\end{proof}

\begin{definition}[Similarity link of an ideal vertex]
\label{thurston-ch03-ideal-link}
\group{ideal-polyhedra}
\source{thurston:page-0052}
\notready
For a hyperbolic manifold assembled from polyhedra, the link \(L(v)\) of an ideal vertex \(v\) is the space of rays through \(v\). It carries a canonical similarity structure modeled on \(E^{n-1}\), obtained by intersecting with horospheres.
\end{definition}

\begin{theorem}[Completeness criterion for ideal polyhedra]
\label{thurston-ch03-ideal-polyhedra-completeness}
\dcref{3.10}
\group{ideal-polyhedra}
\source{thurston:page-0052,thurston:page-0053}
\notready
A hyperbolic manifold obtained by gluing polyhedra with ideal vertices is complete if and only if every ideal-vertex link has Euclidean, rather than merely similarity, holonomy. In dimension three, an oriented link is a torus. If its similarity holonomy is not isometric, it contains a contraction, all holonomy fixes one point, and after translating that point to zero the link is a \((\C^*,\C-0)\)-manifold, equivalently a \((\C,\C)\)-manifold after taking logarithms. Such structures are complete and are determined on a closed oriented surface by a pair \((z_1,z_2)\) linearly independent over \(\R\).
\end{theorem}
\begin{proof}
\sketch The horospherical metric identifies the end with the similarity link, so the ideal end is complete exactly when its similarity holonomy preserves a Euclidean metric. In dimension three, Gauss--Bonnet gives Euler characteristic zero and orientation gives a torus. A nonisometric torus holonomy contains a contraction; commutativity forces every element to fix its unique fixed point, reducing the structure to multiplication on \(\C-0\). The logarithm converts multiplication to translations on \(\C\), and exponentiation gives the converse.
\end{proof}
